\documentclass[12pt]{article}
\usepackage[utf8]{inputenc}
\usepackage[T1]{fontenc}
\usepackage[spanish, es-tabla]{babel}
\usepackage{amsmath}
\usepackage{amsfonts}
\usepackage{amssymb}
\usepackage{graphicx}
\usepackage{geometry}
\usepackage{booktabs}
\usepackage{siunitx}
\usepackage[american]{circuitikz}
\usepackage{enumitem}
\usepackage{caption}
\usepackage{multicol}
\usepackage{makecell}
\usepackage{xcolor}

\geometry{a4paper, margin=0.5in}
\ctikzset{resistor = european}
\sisetup{output-decimal-marker={,}}

\title{\textbf{MAXIMA TRANSFERENCIA DE POTENCIA}}
\author{}
\date{}

\begin{document}

% --- FROM -01.tex ---
\begin{center}
    \large
    PRACTICA 24
    \vspace{0.5cm}

    \huge
    \textbf{MAXIMA TRANSFERENCIA DE POTENCIA}
\end{center}

\section*{FINALIDADES}
\begin{enumerate}
    \item Determinar experimentalmente la condición necesaria para la máxima transferencia de potencia desde una fuente de c.c. hasta una carga.
    \item Comprobar experimentalmente la condición para la máxima transferencia de potencia, indicada en 1.
\end{enumerate}

\section*{INFORMACION PRELIMINAR}

Se ha definido la potencia como velocidad de producción de trabajo. Eléctricamente, la unidad de potencia es el vatio (W). La relación de dependencia entre la potencia de c.c. (W) en una resistencia R, la tensión V entre los extremos de R, y la corriente I en R viene dada por las ecuaciones:

\begin{equation}
W = V \times I = I^2 R = \frac{V^2}{R}
\end{equation}

El suministro de potencia eléctrica a una carga $R_L$ implica una fuente de alimentación V, y una red entre V y la resistencia de carga $R_L$. Consideremos las relaciones de potencia en un circuito sencillo que contiene un generador V de resistencia interna R y que entrega potencia a una resistencia de carga $R_L$.

La corriente I en $R_L$ viene dada por

\begin{equation}
I = \frac{V}{R + R_L}
\end{equation}

Aplicando la fórmula de la potencia $I^2 R$ es evidente que la potencia desarrollada en $R_L$ es

\begin{equation}
P_L = \left(\frac{V}{R + R_L}\right)^2 \times R_L = \frac{V^2 R_L}{(R + R_L)^2}
\end{equation}

Si V es una fuente constante con resistencia interna fija R, ¿con qué valor de $R_L$ habrá la máxima transferencia de potencia desde V hasta $R_L$? Un análisis matemático para hallar este valor requiere el uso del cálculo integral, pero siguiendo un método experimental se podrá determinar el valor de $R_L$.

Supongamos que la tensión V es 100 voltios y que la resistencia R es 100 ohmios. Supongamos también que $R_L$ toma una serie de valores como los indicados en la tabla 24-1. La potencia W, calculada por la fórmula (24-3), está también indicada en la tabla 24-1.

La tabla indica que cuando $R_L$ aumenta de 0 a 100 ohmios, el número de vatios disipados por $R_L$ aumenta desde 0 hasta un máximo de 25. Cuando $R_L$ aumenta de 100 a 100.000 ohmios, el número de vatios transferidos a $R_L$ disminuye desde 25 hasta 0,099.

En este circuito resulta que la máxima transferencia de potencia tiene lugar cuando la resistencia de la carga es igual a la resistencia interna del generador. ¿Es válida esta conclusión para otros generadores V', con resistencia interna R', que transfieren potencia a una carga $R_L$? La respuesta a esta pregunta se puede determinar tomando valores fortuitos de una tensión V', una resistencia R', variando la carga $R_L$ y luego calculando la potencia en $R_L$. Nuevamente llegaremos al resultado de que la máxima potencia se transferirá a la carga cuando $R_L = R'$, resistencia interna del generador.

Podemos pues enunciar la ley que rige la máxima transferencia de potencia a una carga en un circuito de c.c.:

\textit{Un generador transfiere la máxima potencia a una carga cuando la resistencia de ésta es igual a la resistencia interna del generador.}

\section*{RESUMEN}
\begin{enumerate}
    % Merged item from -01 end and -02 start
    \item La potencia W, en vatios, disipada por una resistencia R (ohmios) entre cuyos extremos hay una tensión de c.c. V (voltios), es $W = V^2/R$.
    \item La potencia W, en vatios, disipada por una resistencia (resistor) R (ohmios) en que hay una corriente de c.c. I (amperios) es $W = I^2R$.
    \item Si la tensión c.c. es V (voltios) entre los extremos de un resistor y la corriente es I (amperios), la potencia disipada por la resistencia es $W = V \times I$.
    \item Cuando es entregada potencia por una fuente de alimentación V, cuya resistencia interna es $R_{in}$, a una carga $R_L$, la máxima transferencia de potencia tiene lugar cuando la resistencia de la carga es igual a la resistencia interna de la fuente, es decir, cuando $R_{in} = R_L$.
\end{enumerate}

% --- FROM -02.tex (Table and Autoexamen) ---

\vspace{0.5cm}
\begin{center}
    \textbf{TABLA 24-1}
\end{center}

\begin{center}
\begin{tabular}{l S[table-format=6.0] S[table-format=5.2]}
    \toprule
    \textbf{$R_L$,} & \textbf{$R+R_L$,} & {\textbf{$W = \frac{V^2 R_L}{(R+R_L)^2}$}} \\
    \textbf{ohmios} & \textbf{ohmios} & \textbf{(vatios)} \\
    \midrule
    0 & 100 & 0 \\
    10 & 110 & 8,26 \\
    20 & 120 & 13,9 \\
    30 & 130 & 17,7 \\
    40 & 140 & 20,4 \\
    50 & 150 & 22,2 \\
    60 & 160 & 23,4 \\
    70 & 170 & 24,2 \\
    80 & 180 & 24,7 \\
    90 & 190 & 24,9 \\
    100 & 200 & 25 \\
    110 & 210 & 24,9 \\
    120 & 220 & 24,8 \\
    130 & 230 & 24,6 \\
    140 & 240 & {24,2+} \\
    150 & 250 & {23,9+} \\
    200 & 300 & 22,2 \\
    400 & 500 & 16 \\
    600 & 700 & 12,25 \\
    800 & 900 & 9,87 \\
    1000 & 1100 & 8,26 \\
    10000 & 10100 & 0,98 \\
    100000 & 100100 & 0,099 \\
    \bottomrule
\end{tabular}
\end{center}

\section*{AUTOEXAMEN}
\begin{enumerate}
    \item La corriente en una resistencia de 120 $\Omega$ es 0,1 A. La potencia en vatios disipada por la resistencia es W = \dotfill W.
    \item La tensión entre los extremos de una resistencia es 12 V y la corriente en ella es de 0,05 A. La potencia disipada por la resistencia es W = \dotfill W.
    \item La tensión entre los extremos de una resistencia de 220 $\Omega$ es 16,0 V. La potencia disipada en la resistencia es W = \dotfill W.
    \item Una fuente de alimentación cuya resistencia interna es 25 $\Omega$ entrega potencia a una carga de 50 $\Omega$ conectada entre sus terminales. Si la tensión entregada por la fuente sin carga es 15 V, la potencia disipada por la carga es W = \dotfill W.
    \item Una fuente de alimentación con una resistencia interna de 25 $\Omega$ entrega potencia a una carga resistiva. Si la tensión sin carga en la salida de la fuente es 50 V, la máxima potencia será entregada a la carga cuya resistencia sea $R_L$ = \dotfill $\Omega$.
    \item La potencia entregada por la fuente de la pregunta 5 es W = \dotfill W.
\end{enumerate}

\section*{MATERIALES NECESARIOS}
\textbf{Equipo:} Fuente de c.c. variable regulada; EVM.

\textbf{Resistencias:} ½ W 1.000 $\Omega$ (330, 470, 1.000 y 2.200 $\Omega$ para la pregunta adicional).

\textbf{Diversos:} Caja de décadas de resistencia o potenciómetros de 10.000 $\Omega$ (potenciómetro de 10.000 $\Omega$ para la pregunta adicional); un interruptor bipolar.

\section*{PROCEDIMIENTO}

\begin{multicols}{2}
\begin{enumerate}[label=\textbf{\arabic*.}]
    \item Conectar el circuito de la figura 24-1. V es la alimentación de tensión regulada, R una resistencia de 1.000 ohmios, $R_L$ una caja de décadas de resistencia o un potenciómetro de 10.000 ohmios conectado como reóstato. Para abrir la carga de modo que la resistencia de $R_L$ completamente aislada de V, se pueda ajustar y medir cuando convenga, se utiliza un interruptor bipolar (dos interruptores mecánicos).
\end{enumerate}
    
    % Figura integrada (from -03)
    \begin{center}
        \begin{circuitikz}[scale=0.9, transform shape]
            \draw (0,3) to[V, l_=$10\,\text{V}$, v=$V$] (0,0);
            \draw (0,3) to[R, l=$R$, l_=$1\,\text{k}\Omega$] (3,3);
            \draw (3,3)
                  to[short, -o] (4,3) node[right] {SIB};
            \draw (4,3) to[short, o-*] (4,2);
            \node[right] at (4.2, 2) {A};
            
            \draw (0,0) to[short] (3,0);
            \draw (3,0)
                  to[short, -o] (4,0) node[right] {SIA};
            \draw (4,0) to[short, o-*] (4,1);
            \node[right] at (4.2, 1) {B};
    
            \draw (4,2) to[vR, l=$R_L$, l_=$10\,\text{k}\Omega$] (4,1);
    
            \draw[dashed] (3.5,3.2) -- (3.5,-0.2);
        \end{circuitikz}
        \captionof{figure}{Circuito experimental para determinar la máxima transferencia de potencia a una carga.}
        \label{fig:circuito-24-1}
    \end{center}

    \begin{enumerate}[label=\textbf{\arabic*.}, start=2]
        \item ...mente acoplados). Ajustar la tensión $V = 10$ voltios y mantenerla en este nivel durante el resto de la práctica.
        \item Abrir el interruptor $S_1$. Ajustar el reóstato para que $R_L$ sea 0 ohmios, medidos con un óhmetro. Cerrar $S_1$. Medir con un voltímetro electrónico la tensión $V_L$ entre los extremos de $R_L$ y anotar el resultado en la tabla 24-2.
        \item Abrir el interruptor $S_1$. Ajustar el reóstato de modo que $R_L$ mida 100 ohmios. Cerrar $S_1$. Medir con un EVM la tensión $V_L$ entre los extremos de $R_L$ y anotarla en la tabla 24-2.
        \item Repetir la operación 3 para cada valor de $R_L$ indicado en la tabla.
        \item Por los valores medidos correspondientes de $V_L$ y $R_L$, calcular la potencia de $R_L$ utilizando la fórmula $W = V_L^2/R_L$. Convertir W en milivatios y anotar en la tabla 24-2. Presentar una muestra del cálculo.
        \item Calcular $W_T$ utilizando la fórmula $W_T = V^2/(R + R_L)^2$. Convertir $W_T$ en milivatios y anotar en la tabla 24-2. Presentar una muestra del cálculo.
        \item Dibujar un gráfico de W en función de $R_L$, correspondiendo $R_L$ al eje de abscisas u horizontal, y W al de ordenadas o vertical.
        \item Dibujar un gráfico de $W_T$ en función de $R_L$.
    \end{enumerate}
    
\end{multicols}

\vspace{3cm}

\begin{center}
    \textbf{TABLA 24-2}
\end{center}

\centering
\begin{tabular}{
    S[table-format=5.0]
    S[table-format=5.0]
    S[table-format=1.3]
    S[table-format=2.3]
    S[table-format=2.3]
}
\toprule
{\makecell{\textbf{$R_L$,} \\ \textbf{ohmios}}} &
{\makecell{\textbf{$R+R_L$,} \\ \textbf{ohmios}}} &
{\makecell{\textbf{$V_L$,} \\ \textbf{voltios}}} &
{\makecell{\textbf{$W = \frac{V_L^2}{R_L}$} \\ \textbf{(milivatios)}}} &
{\makecell{\textbf{$W_T = \frac{V^2}{(R+R_L)^2}$} \\ \textbf{(milivatios)}}} \\
\midrule
0 & & & & \\
100 & & & & \\
200 & & & & \\
400 & & & & \\
600 & & & & \\
800 & & & & \\
850 & & & & \\
900 & & & & \\
950 & & & & \\
1000 & & & & \\
1100 & & & & \\
1200 & & & & \\
1500 & & & & \\
1700 & & & & \\
2000 & & & & \\
4000 & & & & \\
6000 & & & & \\
8000 & & & & \\
10000 & & & & \\
\bottomrule
\end{tabular}

% --- FROM -04.tex (Adicional, Preguntas) ---

\subsection*{Adicional}
\begin{enumerate}[start=10, label=\textbf{\arabic*.}]
    \item Describir detalladamente el procedimiento experimental que se utilizaría para determinar el valor de $R_L$ que proporciona la máxima transferencia de potencia en el circuito de la figura 24-2 a $R_L$. Dibujar un esquema del circuito experimental. Tabular las mediciones hechas, incluyendo una columna para $R_L$ y otra para $W$. Cuando se haya determinado experimentalmente el valor de $R_L$, compararlo con el valor calculado de $R_L$, que se halla transformando el circuito de la figura 24-2 en uno equivalente, como el de la figura 24-3.
\end{enumerate}

\textbf{SUGERENCIA:} Conectar el circuito de la figura 24-2, estando el interruptor S abierto y la fuente de alimentación desconectada. Cortocircuitar los puntos C y D. Medir experimentalmente la resistencia entre los puntos A y B. Si el valor medido es $R_T$, la máxima transferencia de potencia tendrá lugar cuando $R_L = R_T$. Elegir un valor de $V$ que dé las lecturas en la región central de la escala del voltímetro cuando se mida $V_L$, tensión entre los extremos de $R_L$.

\begin{figure}[h!]
    \centering
    \begin{circuitikz}[scale=0.8, transform shape]
        \coordinate (origin) at (0,0);
        \coordinate (top_left) at (0,3);
        \coordinate (A_coord) at (2,3);
        \coordinate (B_coord) at (2,0);
        \coordinate (C_coord) at (8,3);
        \coordinate (D_coord) at (8,0);

        \draw (origin) to[V, l_=$V$] (top_left)
              to[opening switch, l=$S$] (A_coord) node[circ, label={above:$A$}] {}
              to[R, l=$R_T$] (C_coord) node[circ, label={above:$C$}] {}
              to[vR, l=$R_L$] (D_coord) node[circ, label={below:$D$}] {}
              to (B_coord) node[circ, label={below:$B$}] {}
              to (origin);
    \end{circuitikz}
    \caption{Circuito equivalente de la fig. 24-2.}
    \label{fig:23-4}
\end{figure}

\subsection*{PREGUNTAS}
\begin{enumerate}[label=\arabic*.]
    \item ¿Para qué valor de $R_L$ es máxima la transferencia de potencia en esta Práctica?
    \item ¿Confirman las mediciones y cálculos de la tabla 24-2 la ley de transferencia de máxima potencia? Explicar cualquier resultado inesperado.
    \item En la figura 24-1, ¿cómo varían las tensiones entre los extremos de $R_L$ con $R_L$? ¿Y la corriente en $R_L$?
    \item En la figura 24-1, ¿cómo varía $W_T$ con $R_L$?
    \item ¿Cuándo se obtiene el mayor porcentaje de disipación de potencia en la carga en el circuito de la figura 24-1? (Ver tabla 24-2 y calcular $W/W_T \times 100$.)
    \item ¿Se obtiene la máxima transferencia de potencia en la figura 24-1 con el mismo valor de $R_L$ para el cual el rendimiento es máximo?
\end{enumerate}

\textbf{NOTA:} El máximo rendimiento tiene lugar cuando se disipa la mayor potencia total ($W_T$) en la carga.

\begin{figure}[h!]
    \centering
    \begin{circuitikz}[scale=0.8, transform shape]
        \node[circ, label={above:$A$}] (A) at (2,4) {};
        \node[circ, label={below:$B$}] (B) at (2,0) {};
        \node[circ, label={above:$C$}] (C) at (11,4) {}; 
        \node[circ, label={below:$D$}] (D) at (11,0) {};
        \node[circ, label={above:$F$}] (F) at (7,4) {};
        \node[circ, label={above:$E$}] (E) at (4,4) {};
        
        \draw (0,4) to[V, l_=$V$] (0,0) to (B);
        \draw (0,4) to[opening switch, l=$S$] (A);
        \draw (B) to[R, l=$470\ \Omega$] (D);
        
        \draw (D) to[vR, l=$R_L$] (C);
        
        % Red de resistores superior con ramas paralelas explícitas
        \draw (A) to[short, -*] (4,4) coordinate (junction_start); 
        \coordinate (junction_end) at (7,4); 
        
        \draw (junction_start) to[R, l=$1000\ \Omega$] (junction_end); 
        \draw (junction_start) -- ++(0,-1.5) to[R, l=$2200\ \Omega$] ++(3,0) -- (junction_end); 
        
        \draw (junction_end) to[R, l=$330\ \Omega$] (C); 
    \end{circuitikz}
    \caption{Determinación del valor de $R_L$ para el cual la transferencia de potencia a $R_L$ es máxima.}
    \label{fig:24-2}
\end{figure}

\subsection*{Respuestas a las preguntas del autoexamen}
Con exactitud de la regla de cálculo.
\begin{enumerate}[label=\arabic*.]
    \item 1,2.
    \item 0,6.
    \item 1,16.
    \item 2,0.
    \item 25.
    \item 25.
\end{enumerate}

\end{document}
